1 & Diseño & Separar la matemática de la interfaz & Una GUI no se puede probar automáticamente & Motor \texttt{nm\_procesar} + funciones \texttt{nm\_*}; la GUI solo llama y grafica & \texttt{pruebas/prueba\_motor.m} corre sin interfaz \\ \midrule
2 & Integración & Simpson 3/8 no era exacto para cúbicas con 7 intervalos & El residuo de 1 intervalo se integraba con trapecio (O(h\textsuperscript{2})) & Residuo de 4 intervalos $\rightarrow$ 2 $\times$ Simpson 1/3 (O(h\textsuperscript{4})) & Exacto para x³ con 3 a 11 intervalos \\ \midrule
3 & Validación EDO & Orden de convergencia "negativo" en Heun y RK4 & Se medía el error solo en t = T, donde 1 − cos(wT) = 0 por coincidencia & Medir el error máximo en toda la trayectoria x(t) & Órdenes 1.11 / 2.00 / 3.94 \\ \midrule
4 & Validación integral & Trapecio parecía más exacto que Simpson & La señal de prueba tenía un número entero de periodos (15 ciclos): el trapecio es superconvergente para funciones periódicas en periodos completos & f = 1.37 Hz (13.7 ciclos) & Simpson 1.4e-7 vs trapecio 1.9e-4 \\ \midrule
5 & Código & \texttt{w} indefinida en \texttt{nm\_senal\_prueba} & Al editar, dos sentencias quedaron en la misma línea detrás de un comentario \texttt{\%} & Separar las líneas & Pruebas pasan \\ \midrule
6 & Validación & Newton local comparado contra \texttt{interp1('pchip')} & pchip es Hermite con preservación de forma, no el mismo polinomio & Marcar "sin equivalente directo" y validar solo contra la solución analítica & Error 1.8e-5 vs analítica \\ \midrule
7 & GUI & Eje derecho de \texttt{yyaxis} persistía al cambiar de vista & \texttt{cla} no elimina el segundo eje Y & \texttt{cla(ax,'reset')} & Revisión de código (pendiente de prueba en MATLAB) \\ \midrule
8 & GUI & Etiquetas de la gráfica de errores sin espacio & \texttt{strcat} elimina espacios finales en argumentos \texttt{char} & Concatenar con \texttt{cellfun} & Revisión de código \\ \midrule
9 & GUI & Posible error al cambiar de archivo con distinto número de ejes & \texttt{Items} e \texttt{ItemsData} del menú con longitudes distintas & \texttt{fijarItems}: vaciar \texttt{ItemsData}, asignar \texttt{Items} y luego \texttt{ItemsData} & Revisión de código \\ \midrule
10 & Física & Velocidad final de 2.09 m/s al integrar la aceleración en reposo & Sesgo y deriva térmica del acelerómetro integrados en el tiempo (drift) & Restar la tendencia por mínimos cuadrados antes de integrar & v(T) $\approx$ 0 m/s \\ \midrule
11 & Mínimos cuadrados & Error al ajustar grado 0 (media) & En la sustitución hacia atrás, \texttt{x(i+1:n)} de un escalar da un vector 1$\times$0 y el producto no es conforme & Sumatoria explícita con ciclo \texttt{for} & MC grado 0 = media (prueba) \\ \midrule
12 & Diseño / física & Restar la tendencia lineal distorsiona el ensayo arranque–frenado del carro & En ese ensayo la aceleración es positiva al inicio y negativa al final; la recta de MC tiene pendiente real que no es drift & 4 opciones de corrección: ninguna, sesgo constante (MC grado 0), tendencia MC, reposo inicial. Para inicio y fin en reposo se usa el sesgo constante & Comparación en \texttt{prueba\_motor.m} \\ \midrule
13 & Datos reales & Los datos del iPhone casi no tienen huecos, entonces la interpolación no se podía evaluar & Con señal real no se conoce el valor verdadero de un dato faltante & \texttt{nm\_simular\_perdida}: elimina en ráfagas un \% de muestras reales y compara el interpolante contra el valor verdadero (RMSE) & Lineal 0.34, Newton 0.12, spline 0.14 m/s\textsuperscript{2} (ejemplo) \\ \midrule
14 & GUI & Riesgo: \texttt{exportgraphics} sobre gráficas de pestañas no visibles & Las pestañas ocultas no se dibujan & Seleccionar la pestaña antes de exportar y restaurarla después & \texttt{prueba\_gui.m} \\ \midrule
15 & Pruebas & Una prueba manual de la GUI no detecta todos los errores & Muchas combinaciones de menús & \texttt{pruebas/prueba\_gui.m}: presiona cada botón, recorre todas las opciones y exporta en 3 formatos sin diálogos & Pendiente de ejecución en MATLAB (Windows) \\ \midrule
16 & Datos reales & phyphox exportó Excel con 3 hojas y el tiempo inicia en −0.0012 s & Formato de exportación del iPhone & \texttt{readtable} lee la primera hoja (Raw Data); el cargador resta t(1). Copias CSV en \texttt{datos/} y originales en \texttt{datos/originales\_phyphox/} & Carga correcta, fs = 100.0 Hz \\ \midrule
17 & Atípicos & El criterio MAD global marcó 106 atípicos en el manejo & La señal no es estacionaria: los baches reales superan el umbral global & Filtro de Hampel local (ventana $\pm$25 muestras, 5$\cdot$1.4826$\cdot$MAD) & 1–5 atípicos por eje; detecta el pico de −29.4 m/s\textsuperscript{2} en t = 24.24 s \\ \midrule
18 & Espectro & La frecuencia dominante del manejo salió en 0.04 Hz & Domina el movimiento lento del carro, no la vibración & Buscar el pico solo en f $\geq$ 5 Hz & 34.2 Hz en z ($\approx$1030 rpm) \\ \midrule
19 & Análisis & Con datos reales la interpolación lineal superó al spline (al revés que en la señal sintética) & $\approx$4 muestras por ciclo a 25 Hz + ruido: los cúbicos oscilan entre muestras & Se documenta: el mejor método depende del muestreo relativo a la frecuencia de la señal & RMSE 0.39 (lineal) vs 0.58 (spline) \\ \midrule
20 & Datos reales & Se esperaba menos vibración en la cabina que en el cofre (aislamiento de los soportes del motor), pero la cabina vibró más (RMS z 0.43 vs 0.29) & La tapa plástica del descansabrazos amplifica la vibración y las rpm eran distintas (motor frío 758 vs caliente 686 rpm) & Hipótesis refutada; se documenta como limitación. Para medir el aislamiento habría que medir los dos puntos a la vez y con la misma rpm & Validación cruzada de frecuencia con phyphox: 22.86 vs 22.65 Hz \\ \midrule
21 & GUI (Mac, tema oscuro) & El botón Exportar, la barra de estado y la curva analítica casi no se veían & Colores fijos pensados para tema claro (fondo claro sin color de texto, línea negra \texttt{k:}) en MATLAB con tema oscuro & Se quitaron los fondos fijos; las líneas usan \texttt{SeriesIndex} (se adapta al tema) y la analítica usa un gris visible en ambos temas & Capturas en Mac R2025 (tema oscuro) \\ \midrule
22 & Validación & Con datos reales la comparación de interpolación salía "0 pts", y la columna "Dif MATLAB" medía otra cosa & Los datos reales no tienen faltantes; la diferencia se calculaba sobre una señal sin huecos & Si no hay faltantes se simula automáticamente un 5 \% y todo se compara en las muestras eliminadas & Cofre: lineal 0.39, Newton 0.53, spline 0.58 m/s\textsuperscript{2}; dif vs interp1 $\approx$ 1e-16 \\ \midrule
23 & Espectro & Con la señal de prueba la frecuencia dominante salía en 5.09 Hz en lugar de 1.37 Hz & El límite de búsqueda ($\geq$ 5 Hz) estaba fijo & Control "f mín. espectro" (5 Hz para datos reales, 0.5 Hz para la señal de prueba) & 1.37 Hz con la señal de prueba \\ \midrule
24 & Figuras & Se necesitan las gráficas en tema claro y oscuro para el documento y la presentación & — & La app exporta cada gráfica en ambos temas (\texttt{theme}, R2025a o posterior). \texttt{generar\_figuras.m} crea 11 figuras con paletas propias en \texttt{figuras/claro} y \texttt{figuras/oscuro}, igual en cualquier versión & Ejecutado en Octave (render de prueba) \\ \midrule
25 & Figuras & Las curvas de convergencia "adelante" y "atrás" se encimaban & Ambas fórmulas tienen exactamente el mismo error máximo para sin t & Estilos distintos (línea continua vs discontinua con marcador más grande) & Inspección visual \\ \midrule
26 & GUI & Al exportar, la app parpadeaba entre tema claro y oscuro y al final quedaba en oscuro & Para exportar en oscuro se cambiaba el tema de toda la ventana (\texttt{theme(app.UIFigure, ...)}) en cada exportación & Cada gráfica se copia (\texttt{copyobj}) a una ventana invisible con el tema deseado y se exporta ahí; la app no cambia. Se agregó el menú "Tema" (automático/claro/oscuro) & \texttt{prueba\_gui.m} en Mac \\ \midrule
27 & Figuras & En la figura de tendencias las líneas casi no se veían & La vibración cruda ($\pm$3 m/s\textsuperscript{2}) tapa tendencias de \textasciitilde{}0.05 m/s\textsuperscript{2} & Se grafica la media móvil de 1 s, eje $\pm$1 m/s\textsuperscript{2} y R\textsuperscript{2} en la leyenda & Revisión de la figura en MATLAB \\ \midrule